
\assignq{2}{3}{BB84: preparation, sifting, intercept--resend and key distillation}{3--4}
\subsection*{(a) Encode bit and basis}
The agreed logical encoding is
\[
\begin{array}{c|cc}
\text{bit}&Z\ \text{basis}&X\ \text{basis}\\\hline
0&\ket H&\ket D=(\ket H+\ket V)/\sqrt2\\
1&\ket V&\ket A=(\ket H-\ket V)/\sqrt2
\end{array}
\]
These two bases are mutually unbiased:
for example \(|\braket{H}{D}|^2=1/2\).

\subsection*{(b) Sift the key and account for loss}
Alice and Bob choose each basis with probability \(1/2\).
The probability of matching choices is
\[
P(\text{match})=P(ZZ)+P(XX)=\tfrac14+\tfrac14
=\boxed{\tfrac12}.
\]
A single photon is detected with probability
\(\eta_{\mathrm{ch}}\eta_d\), assuming no
dark counts or other imperfections:
\[
\boxed{\mathbb{E}[N_{\mathrm{det}}]
=N\eta_{\mathrm{ch}}\eta_d},\qquad
\boxed{\mathbb{E}[N_{\mathrm{sift}}]
=\tfrac12N\eta_{\mathrm{ch}}\eta_d}.
\]
The sifted-key expectation refers to detected
photons for which their announced bases agree.

\subsection*{(c) Derive the error from intercept--resend}
Compute the quantum bit error rate \emph{conditioned}
on Alice and Bob having used the same basis.
Eve chooses that basis with probability \(1/2\);
she then learns and resends the bit without error.
With probability \(1/2\), Eve chooses the
conjugate basis. Her outcome is random relative
to Alice; when Bob measures in Alice's basis
he obtains the wrong bit with probability \(1/2\).
Thus
\[
\boxed{Q=\tfrac12(0)+\tfrac12(\tfrac12)
=\frac14=25\%.}
\]
This is an ideal single-photon intercept--resend
attack, not a universal error rate for every attack.

\subsection*{(d.i) Probability of seeing an error}
For \(m=20\) independently sampled sifted bits,
\[
\boxed{P(\text{at least one error})
=1-(1-Q)^m=1-\left(\frac34\right)^{20}
=0.99683\ \text{(about }99.68\%\text{)}.}
\]

\subsection*{(d.ii) Why discard public test bits?}
The comparison exposes each tested bit on the
authenticated public channel. Eve can therefore
know their values, so the publicly revealed bits
cannot be part of the secret final key.

\subsection*{(d.iii) The two final processing steps}
\textbf{Error correction} reconciles any
differences between Alice's and Bob's sifted
strings, using public parity/syndrome
information as necessary. This discussion
also leaks some information.
\textbf{Privacy amplification} compresses the
reconciled string with an appropriate hash
function to reduce a potential eavesdropper's
remaining information to a negligible level.
\methodnote{Do not apply the basis-sifting factor
twice. For an intercept--resend QBER always
condition on a basis match between Alice and Bob.}
